\section{Planning: controlador de estrategia en tareas de largo alcance}

\subsection{Los límites del papel de la planificación}
El curso define Planning como «asignar el presupuesto de acciones futuras bajo restricciones», y no como generar una lista de pendientes estática. Una planificación de alta calidad debe permitir la replanificación dinámica (dynamic replanning), porque la retroalimentación del entorno cambia continuamente la secuencia de acciones óptima.

\begin{figure}[H]
\centering
\includegraphics[width=0.84\textwidth]{slide-05-planning-loop.jpg}
\caption{Diapositiva del video, aproximadamente en 01:08:20: la planificación necesita formar un ciclo cerrado con la ejecución y la retroalimentación, en lugar de desplegarse de una sola vez}
\end{figure}

\begin{importantbox}{Los tres elementos necesarios del ciclo cerrado de planificación}
\begin{enumerate}
    \item Descomposición del objetivo: mapear el objetivo final a hitos intermedios verificables;
    \item Monitoreo de la ejecución: comparar continuamente el estado esperado con el estado real;
    \item Reescritura del plan: activar la replanificación cuando la desviación supere el umbral.
\end{enumerate}
\end{importantbox}

\subsection{La combinación de planificación y búsqueda}
El punto de vista de Yu Su es que «la planificación no sustituye a la búsqueda, ni la búsqueda sustituye a la planificación». La planificación determina la estructura del espacio de búsqueda, y la búsqueda determina el óptimo local de la acción. En la ingeniería de sistemas conviene elegir si introducir MCTS, beam search o generación heurística de candidatos según la complejidad de la tarea.

\begin{knowledgebox}{Cuándo se necesita búsqueda explícita}
\begin{itemize}
    \item Cuando las consecuencias de la acción se manifiestan con retraso y son irreversibles;
    \item Cuando las trampas de óptimo local son evidentes;
    \item Cuando la señal del evaluador de un solo paso es débil y ruidosa.
\end{itemize}
\end{knowledgebox}

\begin{warningbox}{El costo de la sobreplanificación}
Una planificación de granularidad demasiado fina ocupa una gran cantidad de contexto y de presupuesto de cómputo, y debilita la capacidad de adaptación en tiempo real durante la fase de ejecución. La granularidad de la planificación debe corresponder a la tasa de cambio del entorno.
\end{warningbox}

\subsection{Pseudocódigo del planificador}
\begin{lstlisting}[caption={Un ciclo simplificado de replanificación, que enfatiza la ``actualización activada por desviación''}]
state = observe()
plan = planner.make_plan(goal, state)

for step in range(budget):
    action = planner.next_action(plan, state)
    result = env.execute(action)
    state = update_state(state, result)
    memory.write(result)

    if planner.need_replan(state, goal):
        plan = planner.replan(goal, state, memory.retrieve(goal))
\end{lstlisting}

\subsection{Resumen de la sección}
Lo esencial de Planning no está en «pensarlo todo con la mayor completitud posible de antemano», sino en «poder corregir con rapidez durante la ejecución». La replanificación dinámica es la línea divisoria entre el éxito y el fracaso de un Agent de largo alcance.

